\subsection{Selection of the Studied Frameworks}
\label{subsec:framework-selection}

After defining the functional scope, we selected the frameworks used to
conduct the coverage evaluation. The objective was not to establish an
exhaustive inventory or ranking of existing frameworks, but to construct a
diverse and reproducible sample covering different approaches to the design
and orchestration of \gls{llm}-based \gls{mas}s.

A framework was eligible for inclusion when it satisfied the following
criteria:

\begin{itemize}
    \item it was publicly accessible and sufficiently documented;
    \item it provided a Python API;
    \item it supported the definition of at least one \gls{llm}-based agent;
    \item it supported the composition of multiple agents, tasks, or workflow
    components;
    \item it provided an identifiable and installable version; and
    \item it was actively maintained at the time of selection.
\end{itemize}

Conversely, we excluded libraries limited to direct \gls{llm} calls, closed
platforms without an accessible API, inactive projects, frameworks exclusively
intended for reinforcement-learning agents, and observability tools that did
not support the construction of \gls{mas}s. Frameworks exclusively
targeting programming languages other than Python were also excluded.

Among the eligible frameworks, we adopted a purposive sampling strategy to
maximize the diversity of the sample. In particular, we sought complementary
approaches to agent specification, task definition, tool use, memory
management, orchestration, control flow, and execution. The selected
frameworks consequently represent systems organized around roles and tasks,
conversational teams, graph-based workflows, typed agent definitions,
pipeline-based architectures, and agent SDKs maintained by different
organizations. The frameworks were selected for this complementarity rather
than for their ability to cover every characteristic retained in the
functional scope; the absence of a characteristic from a framework is itself
an outcome of the subsequent coverage analysis.

The resulting sample comprises CrewAI, Google ADK, LangGraph, OpenAI Agents
SDK, Microsoft Agent Framework, Strands Agents, PydanticAI, AutoGen, Semantic
Kernel, and Haystack. Table~\ref{tab:selected-frameworks} presents the main
paradigm for which each framework was included. This classification indicates
the principal contribution of each framework to the diversity of the sample;
it does not imply that a framework is limited to a single paradigm.

\begin{table}[t]
    \centering
    \caption{Selected frameworks and their main represented paradigms.}
    \label{tab:selected-frameworks}
    \small
    \begin{tabularx}{\columnwidth}{
        @{}
        p{0.28\columnwidth}
        >{\raggedright\arraybackslash}X
        @{}
    }
        \toprule
        \textbf{Framework} & \textbf{Main represented paradigm} \\
        \midrule
        CrewAI
            & Role- and task-based agent organization \\
        Google ADK
            & Hierarchical agents and workflow agents \\
        LangGraph
            & Graph-based orchestration with state management \\
        OpenAI Agents SDK
            & Agents, agent handoffs, and tracing \\
        Microsoft Agent Framework
            & Enterprise workflows and fan-out/fan-in orchestration \\
        Strands Agents
            & Lightweight agent composition and agent graphs \\
        PydanticAI
            & Typed agents and structured outputs \\
        AutoGen
            & Agent conversations and teams \\
        Semantic Kernel
            & Agents, services, and enterprise plugins \\
        Haystack
            & Component pipelines and agent workflows \\
        \bottomrule
    \end{tabularx}
\end{table}

To support traceability and reproducibility, we fixed one reference version
for each framework. These versions correspond to the minimum versions
explicitly targeted by the GEAR connectors at the time of the study. For each
framework, we recorded the reference version, its release date, its official
documentation, and the reason for its inclusion.
Table~\ref{tab:framework-metadata} reports these metadata.

\begin{table*}[t]
    \centering
    \caption{Metadata documenting the selection of the studied frameworks.
    The versions correspond to the reference versions targeted by GEAR.}
    \label{tab:framework-metadata}
    \small
    \begin{tabular}{
        @{}
        p{0.18\textwidth}
        p{0.08\textwidth}
        p{0.10\textwidth}
        p{0.17\textwidth}
        p{0.32\textwidth}
        @{}
    }
        \toprule
        \textbf{Framework}
        & \textbf{Version}
        & \textbf{Release date}
        & \textbf{Official source}
        & \textbf{Reason for inclusion} \\
        \midrule

        CrewAI
        & 1.15.4
        & 2026-07-17
        & \href{https://docs.crewai.com/}{Documentation}
        & Role- and task-based orchestration \\

        Google ADK
        & 2.5.0
        & 2026-07-16
        & \href{https://adk.dev/}{Documentation}
        & Hierarchical agents and graph-based workflow agents \\

        LangGraph
        & 1.0.0
        & 2025-10-17
        & \href{https://docs.langchain.com/oss/python/langgraph/overview}
        {Documentation}
        & Stateful graph-based workflows with branching and joins \\

        OpenAI Agents SDK
        & 0.18.0
        & 2026-07-07
        & \href{https://openai.github.io/openai-agents-python/}
        {Documentation}
        & Agent handoffs and native tracing capabilities \\

        Microsoft Agent Framework
        & 1.11.0
        & 2026-07-10
        & \href{https://learn.microsoft.com/en-us/agent-framework/}
        {Documentation}
        & Enterprise agent workflows with fan-out and synchronized fan-in
        orchestration \\

        Strands Agents
        & 1.48.0
        & 2026-07-17
        & \href{https://strandsagents.com/docs/}
        {Documentation}
        & Lightweight composition of deterministic multi-agent graphs \\

        PydanticAI
        & 2.13.0
        & 2026-07-18
        & \href{https://pydantic.dev/docs/ai/overview/}
        {Documentation}
        & Type-safe agents and structured outputs \\

        AutoGen
        & 0.7.5
        & 2025-09-30
        & \href{https://microsoft.github.io/autogen/stable/user-guide/agentchat-user-guide/index.html}
        {Documentation}
        & Conversational agents and bounded agent teams \\

        Semantic Kernel
        & 1.44.0
        & 2026-07-07
        & \href{https://learn.microsoft.com/en-us/semantic-kernel/frameworks/agent/}
        {Documentation}
        & Service- and plugin-oriented agent architecture \\

        Haystack
        & 2.31.0
        & 2026-07-08
        & \href{https://docs.haystack.deepset.ai/docs/intro}
        {Documentation}
        & Pipeline-oriented agent and component composition \\

        \bottomrule
    \end{tabular}
\end{table*}

The functional scope defines the characteristics examined in the study,
whereas the selected frameworks define the systems against which GEAR is
evaluated. Together, these two elements provide the basis for the
bidirectional functional coverage analysis.